% ============================================================
% Capítulo 4 — Planteamiento
% ============================================================

\chapter{Planteamiento}
\label{cap:planteamiento}

Los capítulos anteriores han presentado las herramientas topológicas
(Capítulo~\ref{cap:fundamentos}) y el problema práctico
(Capítulo~\ref{cap:contexto}). Este capítulo explica cómo se combinan
en el trabajo. La idea guía es no proponer mejoras a ciegas: primero se
\emph{diagnostica} qué errores comete un modelo de referencia
(Sección~\ref{sec:diagnostico}). Después se reformula el post-procesado
en lenguaje topológico, lo que permite leer esos errores en un diagrama
de persistencia (Sección~\ref{sec:lectura-topologica}). De ahí salen las
hipótesis y la estrategia de mejora (Sección~\ref{sec:estrategia}) y el
modo de medir si una mejora funciona (Sección~\ref{sec:medir}).

% ------------------------------------------------------------
\section{Diagnóstico de la baseline}
\label{sec:diagnostico}
% ------------------------------------------------------------

El punto de partida es una baseline estándar: una U-Net entrenada con
BCE + Dice cuyas instancias se obtienen umbralizando el mapa de
probabilidades y tomando componentes conexas. Los detalles se describen
en el Capítulo~\ref{cap:metodologia}. Lo que interesa aquí es
\emph{cómo} se equivoca.

\subsection{Taxonomía de errores}

PQ solo distingue entre aciertos, falsos positivos y falsos negativos.
Para decidir qué mejorar hace falta más detalle, así que clasificamos cada
filamento de referencia y cada predicción según cómo se relacionan.
Decimos que una predicción y un filamento de referencia están
\emph{relacionados} si su IoU es al menos $0{,}1$ o si uno cubre al menos
el $50\,\%$ del otro. El segundo criterio evita que un trozo pequeño de un
filamento largo, cuyo IoU es bajo porque la unión es grande, se cuente
como espurio en lugar de como fragmento. Con esta relación, cada
filamento de referencia no acertado se clasifica en uno de estos tipos
(Tabla~\ref{tab:taxonomia}):

\begin{table}[ht]
\centering
\small
\begin{tabular}{@{}lp{9cm}@{}}
\toprule
Tipo & Descripción\\
\midrule
Perdido & Ninguna predicción relacionada.\\
Fragmentado (1:$n$) & Relacionado con varias predicciones (sus trozos).\\
Fusionado ($n$:1) & Su predicción relacionada lo está también con otros filamentos.\\
Mal delimitado & Una predicción relacionada, pero con $\mathrm{IoU}\le0{,}5$.\\
Espurio (predicción) & Predicción sin ningún filamento relacionado.\\
\bottomrule
\end{tabular}
\caption{Tipos de error considerados en el diagnóstico.}
\label{tab:taxonomia}
\end{table}

Esta relación es más exigente que la de las gráficas oficiales de la
competición, que consideran relacionados una predicción y un filamento en
cuanto se solapan ($\mathrm{IoU}>0$). En el
Capítulo~\ref{cap:resultados} los tipos de error se cuentan con la
definición oficial, para que sean comparables con la rúbrica.

Para los perdidos y los fragmentados medimos además su
\emph{recuperabilidad}, es decir, si el mapa de probabilidades contiene
señal que un post-proceso pueda aprovechar. Para un filamento perdido, se
mira la probabilidad máxima en su interior: si es muy baja, el modelo no
lo vio y ningún post-proceso lo recuperará. Para uno fragmentado, se
calcula su \emph{collado}: el umbral más alto al que sus trozos se unen
dentro del filamento. Un collado de $0$ indica un hueco real en la
predicción, que ningún umbral cierra.

\subsection{Dónde falla el modelo}

Este diagnóstico se hizo con una primera versión de la baseline,
entrenada antes de fijar el protocolo del
Capítulo~\ref{cap:metodologia}. Sus cifras no son comparables con las
del Capítulo~\ref{cap:resultados}, pero sirvieron para decidir la
estrategia. El análisis sobre la validación revela varios patrones:
\begin{itemize}
  \item El \emph{recall} cae en los filamentos pequeños (25\,\%), en los
  cercanos al limbo (15\,\%) y en los que ningún otro anotador marca
  (14\,\%, frente al 68\,\% de los que marcan todos).
  \item Muchos filamentos sí se detectan, pero los elimina el filtro de
  área mínima: 203 filamentos de validación se pierden solo por él.
  \item Los espurios tienen casi la misma confianza que los aciertos
  (probabilidad máxima mediana $0{,}917$ frente a $0{,}933$). La confianza
  por sí sola no los separa.
  \item Aproximadamente el 54\,\% de los espurios los marca otro
  anotador de la misma imagen. Más que errores del modelo, son
  desacuerdos entre anotadores.
\end{itemize}

\subsection{El techo humano}

Comparando entre sí las anotaciones de imágenes con varios anotadores, con
el mismo PQ, se obtiene un acuerdo humano de $\mathrm{PQ}\approx0{,}34$.
Esa primera baseline, comparada con cada anotador, obtiene
$\approx0{,}38$. El modelo se parece más a cada anotador que los
anotadores entre sí, porque aprende una especie de consenso. Este dato
es esencial para interpretar cualquier resultado: con referencias tan
ruidosas, un PQ cercano a $1$ no es alcanzable, y parte de los
«errores» son desacuerdos entre expertos.

% ------------------------------------------------------------
\section{Lectura topológica del post-procesado}
\label{sec:lectura-topologica}
% ------------------------------------------------------------

El post-proceso de la baseline tiene una interpretación topológica
directa. Dado el mapa de probabilidades $p$ y un umbral $t$, toma el
conjunto de supernivel
\[
  X_t=\{x : p(x)>t\}
\]
y declara que cada componente conexa de $X_t$ es un filamento. Las
componentes se calculan con 8-conectividad, es decir, con la
construcción T (Observación~\ref{obs:conectividad}), para no fragmentar
artificialmente las barbas de los filamentos. Así, el número de
instancias predichas es $\beta_0(X_t)$.

En lugar de mirar un solo umbral, podemos considerar la filtración de
superniveles completa, con $t$ bajando de $1$ a $0$, y su diagrama de
persistencia en $H_0$ (Sección~\ref{sec:persistencia}). Cada punto
$(b,d)$ es una componente que \emph{nace} en un máximo local de $p$, con
valor $b$, y \emph{muere} en $d$, al fundirse en un punto de silla con
otra más antigua, según la regla del mayor. Una componente está presente
en $X_t$ si y solo si $b>t\ge d$. Por tanto:

\begin{quote}
\emph{El número de instancias que produce la baseline, antes del filtro
de área, es el número de puntos del diagrama en el cuadrante
$b>t,\ d\le t$.}
\end{quote}

Cada tipo de error deja una firma reconocible en el diagrama
(Tabla~\ref{tab:firmas}). Así, el post-procesado deja de ser un ajuste de
dos números y pasa a ser una decisión sobre qué puntos del diagrama
convertir en instancias.

\begin{table}[ht]
\centering
\small
\begin{tabular}{@{}lp{9.5cm}@{}}
\toprule
Error & Firma en el diagrama $H_0$\\
\midrule
Fragmentado & Varios puntos nacidos dentro del filamento en el cuadrante. Si el hueco es recuperable, sus muertes $d$ están justo por debajo de $t$.\\
Fusionado & Una sola componente cubre varios filamentos: el punto de silla que los separa tiene $d>t$, o no hay máximos separados.\\
Perdido & Ningún punto con $b>t$ dentro del filamento. Si hay uno con $b$ algo menor que $t$, hay señal recuperable.\\
Espurio & Punto aislado en el cuadrante, a menudo con poca persistencia $b-d$.\\
\bottomrule
\end{tabular}
\caption{Firma de cada tipo de error en el diagrama de persistencia $H_0$ de la filtración de superniveles.}
\label{tab:firmas}
\end{table}

\begin{observacion}
En la práctica el diagrama se calcula sobre el \emph{logit}
$z=\log\frac{p}{1-p}$ y no sobre $p$. Como la sigmoide es estrictamente
creciente, $\{p>t\}=\{z>\operatorname{logit}(t)\}$: la filtración es la
misma, reparametrizada, y el diagrama tiene los mismos emparejamientos.
La ventaja es numérica: la red está muy segura y, en precisión reducida,
muchos valores de $p$ se redondean a $1$, con lo que se pierde el orden
entre máximos.
\end{observacion}

% ------------------------------------------------------------
\section{Hipótesis y estrategia}
\label{sec:estrategia}
% ------------------------------------------------------------

Del diagnóstico se deriva una estrategia en la que cada línea de mejora
ataca un tipo de error concreto (Tabla~\ref{tab:estrategia}). Las líneas
se ordenan de menor a mayor coste: primero las que no requieren
reentrenar.

\begin{table}[ht]
\centering
\small
\begin{tabular}{@{}llll@{}}
\toprule
Línea & Error al que apunta & Reentrena & En este trabajo\\
\midrule
E1. Umbral y área mínima & Espurios pequeños & No & Sí\\
E2. Post-proceso por persistencia & Fragmentación, espurios & No & Planteada\\
E3. Pérdidas (IoU, clDice, Topograph) & Fragmentación, barbas, continuidad & Sí & Sí\\
E4. Etiquetas multianotador & Desacuerdo entre anotadores & Sí & Planteada\\
\bottomrule
\end{tabular}
\caption{Líneas de mejora, error al que apunta cada una y si se ha
desarrollado en este trabajo.}
\label{tab:estrategia}
\end{table}

\paragraph{E1. Ajuste del post-proceso clásico.} Se ajustan el umbral $t$
y el área mínima de las componentes. Es la mejora más barata y sirve de
segunda referencia para las demás. En lenguaje topológico, el área mínima
es un filtro sobre las componentes de $X_t$ que ignora su persistencia.

\paragraph{E2. Post-proceso basado en persistencia.} La lectura de la
Sección~\ref{sec:lectura-topologica} sugiere sustituir el corte fijo por
decisiones sobre el diagrama, con dos operaciones:
\begin{itemize}
  \item \emph{Filtro por persistencia:} descartar las componentes con
  persistencia $b-d<\tau$, que son poco robustas al umbral, en lugar de
  descartarlas por su área.
  \item \emph{Fusión por proximidad del collado:} unir en una instancia
  las componentes de $X_t$ que se funden en $X_{t-\delta}$, es decir,
  cuyo punto de silla está a menos de $\delta$ por debajo del umbral. Es
  la operación que corregiría los fragmentados recuperables, a costa de
  arriesgarse a crear fusiones.
\end{itemize}
Los parámetros $(\tau,\delta)$ se ajustarían en validación. El diagrama
$H_0$ se calcula con \emph{union-find} en tiempo casi lineal
(Sección~\ref{sec:persistencia-eficiente}), así que el coste añadido es
despreciable frente a la inferencia.

\paragraph{E3. Pérdidas con sesgo topológico.} El post-proceso solo puede
trabajar con la señal que el modelo ya produce. Para los fragmentados con
hueco real y las barbas perdidas hay que cambiar lo que aprende la red.
Se comparan tres alternativas a Dice:
\begin{itemize}
  \item \emph{IoU}, como control: es otra pérdida de solapamiento, sin
  sesgo topológico.
  \item \emph{clDice}~\cite{shit2021}, que mide el solapamiento de los
  esqueletos. Premia la continuidad de las estructuras alargadas y debería
  afinar la detección de las barbas.
  \item \emph{Topograph}~\cite{lux2025}, que penaliza solo las regiones
  topológicamente críticas, con garantías formales
  (Sección~\ref{sec:topograph}) y coste casi lineal.
\end{itemize}

\paragraph{E4. Estrategias de etiquetas.} Con varias anotaciones por
imagen, se comparan tres formas de construir la etiqueta de
entrenamiento: cada anotación como un ejemplo independiente (la opción de
la baseline), la media de las anotaciones como etiqueta blanda, y los
filamentos marcados por la mayoría de anotadores. La referencia de
evaluación no cambia con la estrategia: siempre se evalúa contra cada
anotación original.

\paragraph{Alcance.} Este trabajo desarrolla por completo las líneas E1
y E3, incluida Topograph, cuyos resultados se presentan en el
Capítulo~\ref{cap:resultados}. Las líneas E2 y E4 quedan planteadas y
se recogen como trabajo futuro en el Capítulo~\ref{cap:conclusiones}.

% ------------------------------------------------------------
\section{Cómo medir una mejora}
\label{sec:medir}
% ------------------------------------------------------------

Una mejora del PQ global puede esconder que se arreglan unos filamentos y
se rompen otros. Por eso cada versión se evalúa de tres formas:
\begin{enumerate}
  \item \textbf{Métrica oficial:} PQ global, con SQ y RQ, y las
  distribuciones de la rúbrica (IoU, Dice y relaciones 1:$n$ y $n$:1).
  \item \textbf{Tipos de error:} cuántos filamentos quedan fragmentados,
  fusionados o sin detectar, y cuántas predicciones son espurias, además
  de una medida específica de las barbas.
  \item \textbf{Separación ajuste--reporte:} todas las decisiones (época,
  umbrales, parámetros del post-proceso) se toman en una mitad de la
  validación y la cifra se reporta en la otra, para no inflar los
  resultados.
\end{enumerate}
El protocolo concreto se detalla en el Capítulo~\ref{cap:metodologia}.
